\section{Методология}
\label{sec:methodology}

\subsection{Формализация задачи и архитектура базовой модели}
Пусть $\mathcal{U}$ — множество пользователей, а $\mathcal{V} = \{1, 2, \dots, |\mathcal{V}|\}$ — замкнутый каталог музыкальных треков мощности $|\mathcal{V}| = 98\,504$. Сессия пользователя $u \in \mathcal{U}$ представляется упорядоченной последовательностью пар «трек-длительность»:
\begin{equation}
    \mathcal{S}_u = \big( (a_1, D_1), (a_2, D_2), \dots, (a_T, D_T) \big), \quad a_t \in \mathcal{V}, \; D_t \in \mathbb{R}^+,
\end{equation}
где $D_t$ — фактическое время непрерывного прослушивания трека $a_t$ в секундах.

В качестве базовой модели используется авторегрессионный причинный декодер Causal GPT-2~\cite{petrov2024aligning}, инициализированный с нуля:
\begin{itemize}
    \item Размер словаря токенов: $|\mathcal{V}| + 1 = 98\,505$ (токены $1 \dots 98\,504$ соответствуют ID треков, индекс $0$ зарезервирован под заполнитель \texttt{<PAD>}).
    \item Размерность эмбеддингов: $d_{\text{model}} = 384$.
    \item Количество блоков внимания: $L = 3$, число голов внимания: $H = 4$.
    \item Контекстное окно истории: $T_{\max} = 64$ токена.
    \item Суммарное число параметров: \textbf{43\,174\,656} (~43.2M), все параметры обучаемы.
\end{itemize}

Обучение базовой модели проводилось на 33.2M взаимодействиях с использованием смешанной точности FP16, оптимизатора AdamW ($\text{lr} = 10^{-3}$, $\text{weight\_decay} = 0.05$), линейного разогрева (10\% шагов) и косинусного затухания скорости обучения при эффективном размере батча 128 (микро-батч 32 $\times$ накопление градиентов 4).

\subsection{Формулировки функций вознаграждения}
Для преодоления слепоты функции потерь к длительности прослушивания мы исследовали три альтернативные математические формулировки скалярного вознаграждения $R(s, a)$ (Рис.~\ref{fig:reward_formulations}):

\begin{figure}[!htbp]
    \centering
    \includegraphics[width=0.86\textwidth]{fig2_reward_formulations_comparison.png}
    \caption{\textbf{Сравнение трёх исследованных формулировок функций вознаграждения.} А: Бинарный порог 30 секунд. В: Непрерывная нормализованная доля дослушивания. С: Композитное вознаграждение со штрафом за повторные треки.}
    \label{fig:reward_formulations}
\end{figure}

\subsubsection{Формулировка A: Бинарный порог скипа}
Базовый сигнал, разделяющий взаимодействия на успешные и скипы:
\begin{equation}
    R_A(s, a) = \begin{cases} 
        +1.0, & \text{если } D(a) \ge 30\,\text{с}, \\ 
        -1.0, & \text{если } D(a) < 30\,\text{с}. 
    \end{cases}
    \label{eq:reward_a}
\end{equation}
Порог 30 секунд является общепринятым индустриальным стандартом стриминга: прослушивания короче 30 секунд однозначно свидетельствуют о негативном опыте пользователя.

\subsubsection{Формулировка B: Доля дослушивания}
Нормализует фактическое время прослушивания $D(a)$ относительно номинальной длительности композиции $T(a)$:
\begin{equation}
    R_B(s, a) = \min\left( \frac{D(a)}{T(a)}, 1.0 \right) \in [0.0, 1.0].
    \label{eq:reward_b}
\end{equation}
Для майнинга пар предпочтений $(y_w \succ y_l)$ вводится минимальный порог разрыва $\Delta R_B = R_B(y_w) - R_B(y_l) \ge 0.40$, исключающий пограничный шум между схожими прослушиваниями.

\subsubsection{Формулировка C: Мультикритериальное композитное вознаграждение}
Разработано для оптимизации в парадигме Actor-Critic RL~\cite{volkov2026session}. Учитывает вероятность скипа, долю дослушивания и штраф за навязчивое повторение недавно прослушанных композиций:
\begin{equation}
    R_C(s, a) = \big(1 - P(\text{skip} \mid s, a)\big) \cdot \text{CompletionRatio}(s, a) - \lambda_{\text{skip}} P(\text{skip} \mid s, a) - \lambda_{\text{rep}} \mathbb{I}(a \in s_{t-3:t}),
    \label{eq:reward_c}
\end{equation}
где $\lambda_{\text{skip}} = 0.5$, $\lambda_{\text{rep}} = 0.3$.

\subsection{Парадигмы выравнивания обучения с подкреплением}

\subsubsection{1. Двухстадийное реранжирование с функцией ценности $Q(s, a)$}
Обучается независимая нейросетевая модель-критик $Q_\phi(s, a)$, оценивающая математическое ожидание вознаграждения $\mathbb{E}[R \mid s, a]$ поверх замороженных эмбеддингов базового трансформера (двухслойный перцептрон, 509K параметров)~\cite{petrov2024aligning}. При выводе итоговая оценка кандидата вычисляется как линейная комбинация логитов политики и оценки критика:
\begin{equation}
    \text{Score}(a \mid s) = z_{\text{base}}(a \mid s) + \lambda \cdot Q_\phi(s, a), \quad \lambda \in \{0.1, 0.3, 0.5\}.
\end{equation}

\subsubsection{2. Actor-Critic PPO с ограничением расхождения}
Классический подход оптимизации суррогатной функции политики~\cite{schulman2017proximal}. Актёр генерирует распределение $\pi_\theta(a \mid s)$, из которого сэмплируются действия. Для предотвращения катастрофического дрейфа оптимизируется целевая функция с ограничением дивергенции Кульбака--Лейблера:
\begin{equation}
    \max_\theta \mathbb{E}_{s \sim \mathcal{D}, a \sim \pi_\theta} \big[ R_C(s, a) \big] - \beta_{\text{KL}} \mathbb{D}_{\text{KL}}(\pi_\theta(\cdot \mid s) \parallel \pi_{\text{ref}}(\cdot \mid s)).
\end{equation}

\subsubsection{3. Прямое согласование предпочтений DPO с якорной регуляризацией}
Вместо аппроксимации нестабильной функции ценности DPO~\cite{rafailov2023direct} аналитически выражает скрытую функцию вознаграждения через отношение вероятностей текущей и опорной политики:
\begin{equation}
    r(s, a) = \beta \log \frac{\pi_\theta(a \mid s)}{\pi_{\text{ref}}(a \mid s)}.
\end{equation}
Обучение проводится на датасете троек $\mathcal{D}_{\text{pref}} = \{(s, y_w, y_l)\}$, где $y_w$ — трек с позитивным прослушиванием ($D \ge 30$ с), а $y_l$ — скип ($D < 30$ с). 

Для сохранения структуры предсказания следующего трека вводится вспомогательная якорная регуляризация правдоподобия:
\begin{equation}
    \mathcal{L}_{\text{total}}(\theta) = \mathcal{L}_{\text{DPO}}(\theta; \pi_{\text{ref}}) + \gamma \cdot \mathcal{L}_{\text{SFT}}(y_w \mid s),
    \label{eq:dpo_sft}
\end{equation}
где
\begin{equation}
    \mathcal{L}_{\text{DPO}}(\theta; \pi_{\text{ref}}) = - \mathbb{E}_{(s, y_w, y_l) \sim \mathcal{D}_{\text{pref}}} \left[ \log \sigma \left( \beta \log \frac{\pi_\theta(y_w \mid s)}{\pi_{\text{ref}}(y_w \mid s)} - \beta \log \frac{\pi_\theta(y_l \mid s)}{\pi_{\text{ref}}(y_l \mid s)} \right) \right],
\end{equation}
а $\mathcal{L}_{\text{SFT}}(y_w \mid s) = - \log \pi_\theta(y_w \mid s)$ максимизирует правдоподобие успешного действия.

\subsection{Протокол валидации и статистический инструментарий}
Оценка проводится на отложенном тестовом множестве из \textbf{9\,048 пользователей} (\texttt{test\_interactions.parquet}) с полной предысторией. Ранжирование осуществляется по \textbf{полному каталогу из 98\,504 треков} без искусственного случайного негативного сэмплирования.

Применялся комплексный статистический аппарат:
\begin{enumerate}
    \item \textbf{Непараметрический бутстрэп:} $B = 5\,000$ ресэмплов с вычислением 95\% доверительных интервалов для метрик HR@K, NDCG@K, MRR и разностей $\Delta = \text{Metric}_{\text{DPO}} - \text{Metric}_{\text{Base}}$.
    \item \textbf{Парный тест МакНемара с поправкой Эдвардса на непрерывность:}
    \begin{equation}
        \chi^2 = \frac{(|n_{01} - n_{10}| - 1)^2}{n_{01} + n_{10}},
        \label{eq:mcnemar}
    \end{equation}
    где $n_{01}$ — число пользователей, угаданных DPO, но неугаданных Base; $n_{10}$ — наоборот.
    \item \textbf{Критерий знаковых рангов Уилкоксона:} Непараметрическая парная проверка гипотезы сдвига распределений ранжирования.
    \item \textbf{Метрики здоровья и разнообразия каталога:}
    \begin{itemize}
        \item \textit{Покрытие каталога (Coverage):} Доля треков каталога, рекомендованных хотя бы раз в Top-10.
        \item \textit{Энтропия Шеннона:} $H = - \sum_{i=1}^{|\mathcal{V}|} p(i) \log_2 p(i)$, отражающая равномерность распределения показов.
        \item \textit{Индекс Джини:} Коэффициент неравенства показов треков.
        \item \textit{Корреляция Спирмена с популярностью:} $\rho_s$, отражающая склонность модели к тиражированию топ-хитов.
    \end{itemize}
    \item \textbf{Мультишаговые авторегрессионные сессии:} Симуляция формирования очереди из 5 треков ($s_t \to \hat{a}_{t+1} \to \dots \to \hat{a}_{t+5}$) без отклика пользователя с измерением пошаговой вероятности скипа $P(\text{skip} \mid s_k, \hat{a}_k)$ и ожидаемой длительности сессии~\cite{volkov2026session}.
\end{enumerate}
